% Scope: Derive spatial phase accumulation and trajectory geometry, with sampling definitions separated from effective resolution.
% Status: architecture only; no chapter prose or completed problems yet.
% Prerequisites: Chapters 1, 3 and 4.
% Planned worked examples: Slice thickness, gradient strength, FOV, k-space extent and pixel spacing.
% Planned figures: Slice selection, Cartesian, radial, spiral and EPI trajectories.
% Planned challenge problems: Design a constrained Cartesian readout; expose a zero-filling resolution fallacy.
\chapter{Spatial Encoding and k-Space}
\label{ch:05}

\section{Gradient fields and spatial phase}
\label{sec:05:gradient-fields-and-spatial-phase}
\subsection{Linear field expansion and gradient units}
\subsection{Phase accumulation and trajectory origin}
\subsection{Refocusing pulses and effective gradient signs}

\section{Slice selection and partition encoding}
\label{sec:05:slice-selection-and-partition-encoding}
\subsection{Frequency-to-position mapping}
\subsection{Slice rephasing and finite pulse duration}
\subsection{Two-dimensional slices versus three-dimensional slabs}

\section{Frequency and phase encoding}
\label{sec:05:frequency-and-phase-encoding}
\subsection{Readout gradients and dwell time}
\subsection{Phase-encoding steps and Cartesian grids}
\subsection{Echo position and prephasing}

\section{Sampling geometry}
\label{sec:05:sampling-geometry}
\subsection{Field of view and k-space spacing}
\subsection{Finite extent, nominal resolution and voxel grid}
\subsection{Nyquist sampling, aliasing and zero filling}

\section{Trajectory families and ordering}
\label{sec:05:trajectory-families-and-ordering}
\subsection{Cartesian and echo-planar trajectories}
\subsection{Radial and spiral sampling}
\subsection{Linear, centric and interleaved view ordering}

\section{Gradient moments and trajectory fidelity}
\label{sec:05:gradient-moments-and-trajectory-fidelity}
\subsection{Zeroth and higher moments}
\subsection{Gradient delays and eddy-current errors}
\subsection{Concomitant fields and nonlinear encoding}

\section{Worked examples}
\label{sec:05:worked-examples}
% Planned calculations: Slice thickness, gradient strength, FOV, k-space extent and pixel spacing.
\subsection{Derivation and quantitative calculation}
\subsection{Assumptions, limiting cases and scanner interpretation}

\section{Challenge problems}
\label{sec:05:challenge-problems}
% Problem design: Design a constrained Cartesian readout; expose a zero-filling resolution fallacy.
% Use challengeproblem and stable labels prob:05:<topic>.
% Complete solutions belong in Appendix E, section for this chapter.
% Use \begin{solution}{prob:05:<topic>} ... \end{solution} there.
\subsection{Derivation and quantitative design}
\subsection{Model criticism and integrated reasoning}
